\section*{5. Complete cascade and simulation results}
The outer position controller and nonlinear inverse generate desired rotor speeds; two IP loops with anti-windup generate the limited motor voltages. With $x=[q^\mathsf{T},\dot q^\mathsf{T},\omega^\mathsf{T}]^\mathsf{T}$, the plant is
\begin{equation}
\boxed{\dot x=\begin{bmatrix}
\dot q\\M_q^{-1}[\Gamma s(\omega)-c_q-g_q]\\-a_m\omega
\end{bmatrix}
+\begin{bmatrix}0_{2\times2}\\0_{2\times2}\\b_mI_2\end{bmatrix}v},
\qquad y=\begin{bmatrix}I_2&0_{2\times4}\end{bmatrix}x.
\end{equation}
Two motor integrators augment the plant, with gains retained from the preceding designs. The tachometer filters $H_\omega(s)=180/(s+180)$ remain in both motor feedback paths: the IP laws of Section~2 use $\omega_f=H_\omega\omega_{\rm meas}$ in place of $\omega$. Reference differentiation follows Section~4. Finite rotor tracking introduces the acceleration discrepancy
\begin{equation}
\ddot q-a=M_q^{-1}\Gamma\bigl[s(\omega)-s(\omega_d)\bigr].
\end{equation}
All runs start from rest and use forward Euler integration with step 0.002 s. The table reports the supplied 60 s yaw/pitch trajectories with sensor effects enabled. The two PD cases use identical gains without/with reference-velocity feedforward (FF); the weighted-integral design is given below.
\begin{center}
\begin{tabular}{llccc}
\hline
Controller & Sensors & Yaw RMS (rad) & Pitch RMS (rad) & Peak $|v|$ (V)\\\hline
Ideal motors & Ideal & 0.09964 & 0.05981 & ---\\
Full IP cascade & Ideal & 0.09941 & 0.05975 & 6.46\\
Full IP cascade & Course & 0.09945 & 0.05969 & 6.91\\
Integral cascade & Course & 0.09943 & 0.05970 & 7.67\\
\hline
\end{tabular}
\end{center}

Errors use true plant angles. Acceleration feedforward is absent, so reference acceleration still forces the error. Separate signed single-axis and simultaneous step tests use ideal measurements and disable reference differentiation: settling is 3.1980--3.2040 s without overshoot. The full cascade is slightly slower than the ideal 3 s design.
\begin{center}
\includegraphics[width=\linewidth]{report_part2_course_baseline.png}
\reportcaption{Course-reference tracking with FF (left), and matched tracking-error and motor-voltage comparisons without/with FF (middle/right).}
\end{center}

\subsection*{Optional: position integral action}
Add an acceleration-valued integral contribution $J$. Matching the ideal per-axis denominator to $(s+\lambda)^3$ gives
\begin{equation}
\dot J=K_i^qe,\qquad a=K_p^q(bq_d-q)+K_d^q(\dot q_d-\dot q)+J,
\quad K_p^q=3\lambda^2I_2,\quad K_d^q=3\lambda I_2,\quad K_i^q=\lambda^3I_2.
\end{equation}
For ideal motors, zero initial conditions and reference differentiation disabled, each axis has
\[
\frac{q_i(s)}{q_{d,i}(s)}=\frac{3b\lambda^2s+\lambda^3}{(s+\lambda)^3}
\quad\xrightarrow{\ b=1/3\ }\quad\frac{\lambda^2}{(s+\lambda)^2}.
\]
Thus $b=1/3$ cancels one stable pole and preserves the ideal critically damped step response. With $\lambda=1.9446$ s$^{-1}$, the same motor filters and ideal step-test measurements, settling is 3.0280--3.0860 s with at most 0.0283\% overshoot. Simultaneous steps reach 24 V; both outer integrators pause during saturation. On the moving course reference, the weighted design has roughly tenfold the RMS error of PD with FF. Preserving the step response through reference weighting does not preserve moving-reference tracking performance.
